\section{What the code does with a pre-generated pair}
\label{sec:code-reading}
The explanation below was checked against the source, not inferred from the results. File and line numbers refer to the repository at the commit that contains this report.

\begin{enumerate}
  \item \textbf{When a dynamic preparation starts.} The planner gives each request a generation layer. For the dynamic policy the candidate layers are tried from the one just
    before the request's own layer, backwards (\code{compiler_trace.py:254}), and the first one with room is taken. The controller launches a layer's preparations when that layer is
    activated (\code{request_scheduler.py:151--158}, \code{compiler_scheduler.py:66}). So a pair for a request in layer $L$ is normally started at the \emph{start} of layer $L-1$
    and has to wait for all of layer $L-1$ to finish before its request can even be released. The lookahead of eight layers is only an upper limit; varying it made no difference
    (Section~\ref{sec:ablations}).
  \item \textbf{What a prepared pair occupies.} A preparation books one memory card at each of its two cores, with a window of 1\,000\,ms
    (\code{adaptive_continuous.py:192--236}), and it may not exceed the compiler cap, the same counter that limits the autonomous protocol
    (\code{adaptive_continuous.py:201}). The booking is released only when the pair is consumed (\code{reservation.py:473}) or its window ends.
    If the cap is reached the preparation is rejected and retried every millisecond until its request starts (\code{compiler_scheduler.py:113--120}).
  \item \textbf{A request that has a ready pair still needs a memory of its own.} Admission asks for a free card at both cores, and a card with any overlapping booking
    is not free (\code{request_scheduler.py:202--215, 232--255}); the prepared pair's own booking overlaps. When the request's reservation runs, the generation protocol
    looks for a ready pair for the same two cores, and if it finds one it \emph{swaps} the pair's state into the request's own memory after a classical delay
    (\code{generation.py:265--296, 301--343}) and only then releases the prepared pair's memory (\code{generation.py:342}). So for a short time a hit holds \emph{two} cards
    at each core. This is a property of ACE's design, not of the compiler policy, and it is why requests served from a ready pair can still wait for admission.
  \item \textbf{What an on-demand request occupies.} One card per core, from admission until completion. Its booking window is the request duration (50\,ms), but the booking is
    released at completion, so what matters is the time it exists, about 6\,ms in the measurements, not the window.
  \item \textbf{CGP and ACGP.} Every cycle of the autonomous protocol draws a neighbour by roulette wheel over a probability table (\code{adaptive_continuous.py:262--276}) and books one card at both
    ends until a 0.2\,s boundary within the next 0.2\,s (\code{adaptive_continuous.py:147--178}; the period is $2\times10^{11}$\,ps, \code{node.py:57}), as long as fewer than the cap are in use.
    The table is uniform for CGP (\code{adaptive_continuous.py:243--260}). For ACGP it is updated every 0.2\,s: each neighbour that appears in an entanglement path used in the last 0.2\,s gains $0.05$
    and the table is renormalized (\code{adaptive_continuous.py:404--447}). Pairs are therefore made for whichever neighbour was drawn, not for a particular request.
    An application request takes a pair for its own two cores if one exists: first a compiler pair targeted at \emph{that} request, otherwise an ordinary autonomous pair with no
    compiler metadata (\code{adaptive_continuous.py:511--523}), through the same swap. A pair nobody uses before its window ends is discarded; that is the waste the experiments count.
\end{enumerate}
Two consequences follow directly from the code and are tested in the measurements below. First, pre-generation always costs memory cards for longer than an on-demand request does
(a ready pair exists \emph{before} its request, and for a later hop of a chain well before it, because a hop is released only after its predecessor completed). Second, it
saves only the generation time (a request served from a ready pair generates in about a quarter of a millisecond instead of about a millisecond); the lead and the classical delays are the same either way.
